\documentclass[11pt]{article}
\usepackage[utf8]{inputenc}
\usepackage[margin=1in]{geometry}
\usepackage{amsmath}
\usepackage{graphicx}
\usepackage{booktabs}
\usepackage{hyperref}
\usepackage[round]{natbib}

\title{Clustered Synapses Emerge in Distinct Dendritic Domains of Visual Cortex Pyramidal Neurons Before Eye Opening}
\author{Author A,$^{1}$ Author B,$^{1}$ Author C,$^{2}$ Author D$^{1,*}$\\
\small $^{1}$Department of Neuroscience, Institute of Neuroscience\\
\small $^{2}$Department of Physiology, University of Neuroscience\\
\small $^{*}$Corresponding author: authord@institute.edu}
\date{}

\begin{document}
\maketitle

\begin{abstract}
Synaptic clustering---the spatial grouping of functionally related synaptic inputs along dendrites---is a hallmark of mature cortical circuits that supports nonlinear dendritic computation. However, when and how this organization emerges during development remains unknown. Here we used concurrent in vivo two-photon calcium imaging and whole-cell patch-clamp recording to map 354 functional synapses producing 3,440 transmission events across 12 dendritic areas from 11 mice aged postnatal day 8 to 13 (P8--P13), before eye opening in primary visual cortex. We demonstrate that functional synapse density increases several-fold during the second postnatal week, with both spine and shaft synapses contributing. At P8--P10, synapses assembled in confined dendritic segments while other segments remained devoid of inputs; by P12--P13, dendrites were nearly entirely covered by domains of co-active synapses. We identified distinct dendritic domains of clustered, co-active synapses that expand and increase in synapse number during development. Critically, local co-activity with neighboring synapses correlated positively with synaptic stabilization and potentiation, while low co-activity predicted synaptic depression. These findings reveal that functionally organized synaptic clustering on layer 2/3 pyramidal cell dendrites is established by spontaneous network activity before the onset of visual experience, through a local co-activity-dependent plasticity mechanism.
\end{abstract}

\section{Introduction}

Dendritic computation relies on the precise spatial arrangement of synaptic inputs along the dendritic arbor. In the adult cortex, neighboring synapses on the same dendritic branch often share similar functional properties---a phenomenon termed synaptic clustering \citep{chen2011clustered, jia2010dendritic, wilson2016orientation}. This spatial organization enables nonlinear summation of co-active inputs within dendritic compartments, greatly expanding the computational capacity of individual neurons \citep{larkum2009synaptic}. Despite the established importance of synaptic clustering in adult circuits, a fundamental question remains: when does this organization first emerge during development, and what mechanisms drive its formation?

Electron microscopy studies have shown that excitatory synapses in the developing cortex appear during the first postnatal week and reach peak density around postnatal day 14 (P14) in mice, coinciding with eye opening \citep{lendvai2000experience}. However, these structural data do not reveal the functional organization of synaptic inputs at the single-synapse level. Previous in vivo calcium imaging studies mapped synaptic inputs in adult cortex \citep{jia2010dendritic, wilson2016orientation}, but no study has achieved this in the neonatal cortex, where many synapses form directly on the dendritic shaft rather than on spines, presenting unique technical challenges.

A critical open question concerns the role of sensory experience versus spontaneous activity in establishing synaptic organization. Before eye opening, the developing visual cortex is driven by spontaneous retinal waves and intrinsic network activity \citep{blankenship2010mechanisms}, which generate coordinated bursts of activity across neuronal populations \citep{valeeva2019emergence}. These spontaneous activity patterns have been proposed to instruct early circuit refinement through Hebbian and spike-timing-dependent plasticity mechanisms \citep{hebb1949organization, markram1997regulation}. However, whether spontaneous activity is sufficient to establish the clustered synaptic organization observed in the adult cortex has not been tested.

A further challenge is the detection of shaft synapses in vivo. Unlike spine synapses, where calcium transients can be readily detected as localized signals on morphologically identifiable protrusions, shaft synapses produce calcium signals directly on the dendrite that can be obscured by back-propagating action potentials. Previous studies were therefore limited to spine synapses in mature tissue. Detecting shaft synapses---which are abundant in developing dendrites before spine maturation---requires suppression of action potentials in the recorded neuron.

Here we address these questions by combining in vivo two-photon calcium imaging with somatic whole-cell voltage-clamp recording in layer 2/3 pyramidal neurons of mouse primary visual cortex between P8 and P13. By holding neurons at $-30$~mV to enhance NMDA receptor-mediated calcium transients and blocking action potentials with intracellular QX-314, we achieved unambiguous detection of individual synaptic transmission events at both spine and shaft synapses \citep{kleindienst2011activity, winnubst2015spontaneous}. Our results reveal that functionally clustered synaptic domains emerge before eye opening through a co-activity-dependent plasticity mechanism driven by spontaneous network activity.

\section{Methods}

\subsection{Animals and In Utero Electroporation}

Timed-pregnant C57BL/6J mice underwent in utero electroporation at embryonic day 16.5 (E16.5) to express GCaMP6s \citep{chen2012ultrasensitive} and DsRed Express in layer 2/3 pyramidal neurons of primary visual cortex. Plasmid concentrations were 2~mg/mL for pCAGGS-GCaMP6s and 0.5--2~mg/mL for pCAGGS-DsRed Express. Electroporation was performed with 5 square-wave pulses at 30~V, 50~ms duration, and 950~ms interval.

\subsection{In Vivo Two-Photon Imaging and Electrophysiology}

Experiments were performed on P8--P13 pups (11 mice total). Anesthesia was induced with 2\% isoflurane and maintained at 0.7--1\% during recording, a concentration that preserves basic spontaneous network activity patterns. A craniotomy was performed over V1, and a glass patch pipette (4.5--6~M$\Omega$) coated with BSA-Alexa 594 was advanced to a labeled neuron under two-photon guidance.

Whole-cell voltage-clamp recordings were obtained using a Multiclamp 700b amplifier with 10~kHz sampling and 3~kHz filtering. The intracellular solution contained (in mM): CsMeSO$_3$ 120, NaCl 8, CsCl$_2$ 15, TEA-Cl 10, HEPES 10, QX-314 bromide 5, MgATP 4, and Na-GTP 0.3. Neurons were held at $-30$~mV to enhance NMDA receptor activation and thereby increase synaptic calcium transients, while QX-314 blocked action potentials to prevent contamination from back-propagating spikes.

\begin{table}[ht]
\centering
\caption{Intracellular Solution Composition}
\label{tab:ics}
\begin{tabular}{lc}
\toprule
\textbf{Component} & \textbf{Concentration (mM)} \\
\midrule
CsMeSO$_3$ & 120 \\
NaCl & 8 \\
CsCl$_2$ & 15 \\
TEA-Cl & 10 \\
HEPES & 10 \\
QX-314 bromide & 5 \\
MgATP & 4 \\
Na-GTP & 0.3 \\
\bottomrule
\end{tabular}
\end{table}

Two-photon imaging was performed with resonant scanning combined with piezo-driven z-positioning for multi-plane acquisition. The DsRed channel was used for structural identification of spines versus shaft, while the GCaMP6s channel detected synaptic calcium transients.

\subsection{Synapse Classification and Analysis}

Synapses were classified as spine or shaft based on morphological analysis of the DsRed structural channel. Spine synapses were identified as calcium transients localized to morphologically identified spine heads. Shaft synapses were identified as calcium transients occurring directly on the dendritic shaft without an associated spine structure; these could be detected only when action potentials were blocked with QX-314.

Synaptic domains were identified as contiguous dendritic segments containing co-active synapses firing synchronously above chance levels. Co-activity scores were computed for each synapse, quantifying the degree of synchronous firing with neighboring synapses within a domain. Changes in transmission frequency over the recording period were used to assess synaptic potentiation and depression, and these changes were correlated with the co-activity score.

\subsection{Statistical Analysis}

Age-dependent changes in synapse density and transmission frequency were assessed using linear regression. Inter-synapse distance distributions were compared between age groups using cumulative distribution analysis. The correlation between co-activity score and change in transmission frequency was assessed with correlation analysis. Domain identification employed a spatial clustering algorithm based on co-activity thresholds.

\section{Results}

\subsection{Mapping Functional Synapses In Vivo During Early Postnatal Development}

We successfully mapped 354 functional synapses producing 3,440 transmission events across 12 dendritic areas from 11 mice between P8 and P13 (Table~\ref{tab:overview}). Both spine and shaft synapses were detected, with shaft synapses being visible only because action potentials were suppressed by intracellular QX-314. Synaptic transmission events could be identified as localized calcium transients at specific dendritic locations, and individual synapses could be distinguished even during spontaneous network barrages when multiple synapses were simultaneously active.

\begin{table}[ht]
\centering
\caption{Overview of Experimental Dataset}
\label{tab:overview}
\begin{tabular}{lc}
\toprule
\textbf{Parameter} & \textbf{Value} \\
\midrule
Number of animals & 11 \\
Age range & P8--P13 \\
Number of dendritic areas & 12 \\
Total functional synapses & 354 \\
Total transmission events & 3,440 \\
Anesthesia & Isoflurane (0.7--1\%) \\
Holding potential & $-30$~mV \\
AP blocker & QX-314 (intracellular) \\
\bottomrule
\end{tabular}
\end{table}

Most synaptic transmission events occurred during barrages of spontaneous network activity, visible as compound synaptic currents at the soma. The somatic whole-cell recording confirmed that imaging-detected calcium transients corresponded to actual synaptic currents, validating the approach for detecting individual synaptic inputs in the developing cortex.

\subsection{Functional Synapse Density Increases Several-Fold During the Second Postnatal Week}

Quantification of functional synapse density revealed a dramatic increase across the P8--P13 developmental window (Table~\ref{tab:density}). At P8--P9, dendrites carried sparse functional inputs with large gaps between neighboring synapses. By P12--P13, synapse density had increased several-fold, with dendrites approaching near-complete coverage by functional synaptic inputs. Linear regression confirmed a significant positive correlation between age and synapse density per 100~$\mu$m of dendrite.

\begin{table}[ht]
\centering
\caption{Developmental Changes in Synapse Density and Transmission}
\label{tab:density}
\begin{tabular}{lccc}
\toprule
\textbf{Age Group} & \textbf{Synapse Density} & \textbf{Transmission Frequency} & \textbf{Inter-Synapse Distance} \\
\midrule
P8--P9 & Low (sparse) & Lower & Larger (shifted rightward) \\
P10--P11 & Intermediate & Intermediate & Intermediate \\
P12--P13 & High (dense) & Higher (several-fold increase) & Smaller (shifted leftward) \\
\bottomrule
\end{tabular}
\end{table}

The frequency of transmission events also increased several-fold from P8 to P13, with both spine and shaft synapses contributing to this increase. Approximately 20\% of synapses were highly active, while the remaining 80\% showed rare transmission events, producing a highly skewed, long-tailed distribution of transmission frequencies. This distribution likely reflects differences in release probability between presynaptic terminals rather than differential firing rates of presynaptic neurons \citep{katz1966release}.

\subsection{Synapses Assemble in Confined Dendritic Segments at Early Ages}

Analysis of the spatial distribution of functional synapses along dendrites revealed a striking organizational principle. At P8--P10, synapses were not uniformly distributed; instead, they assembled in confined dendritic segments while other segments remained entirely devoid of functional inputs. Cumulative distribution analysis of inter-synapse distances confirmed that younger dendrites (P8--P10) had significantly larger spacing between neighboring synapses compared to older dendrites (P11--P13).

By P12--P13, the patchy distribution had resolved into nearly continuous synaptic coverage, with most dendritic segments carrying functional inputs. This transition from confined, clustered input zones to complete dendritic coverage occurred over just a few days, suggesting rapid synapse addition during this critical developmental window.

\subsection{Distinct Dendritic Domains of Co-Active Synapses}

To determine whether the spatial clustering of synapses reflected functional organization, we analyzed the co-activity patterns of neighboring synapses and identified distinct dendritic domains---contiguous segments containing synapses that fired synchronously above chance levels (Table~\ref{tab:domains}).

\begin{table}[ht]
\centering
\caption{Dendritic Domain Properties Over Development}
\label{tab:domains}
\begin{tabular}{lccl}
\toprule
\textbf{Property} & \textbf{P8--P10} & \textbf{P12--P13} & \textbf{Trend} \\
\midrule
Domain extent ($\mu$m) & Smaller & Larger & Increases with age \\
Synapses per domain & Fewer & More & Increases with age \\
Domains per dendrite & Few & More & Increases with age \\
Dendrite coverage & Low & High ($\sim$entire) & Dramatic increase \\
\bottomrule
\end{tabular}
\end{table}

At P8--P10, individual domains were small, containing only a few co-active synapses within short dendritic stretches. By P12--P13, domains had expanded considerably: they covered longer dendritic segments, contained more synapses, and collectively covered nearly the entire dendritic arbor. High-activity synapses (the approximately 20\% with frequent transmission) were found within these domains, further supporting the interpretation that domain organization reflects functionally meaningful synaptic clustering.

\subsection{Local Co-Activity Predicts Synaptic Fate}

We next tested whether local co-activity with neighboring synapses influences the fate of individual synapses over the course of the recording. We computed a co-activity score for each synapse and correlated it with the change in transmission frequency during the recording session.

Synapses with high co-activity scores---those firing synchronously with their dendritic neighbors---tended to increase their transmission frequency (potentiation/stabilization). Conversely, synapses with low co-activity scores---those firing out of synchrony with their neighbors---tended to decrease their transmission frequency or cease transmitting altogether (depression/elimination). Correlation analysis confirmed a significant positive relationship between co-activity score and change in transmission frequency, indicating that local synchrony with neighbors promotes synaptic potentiation.

This co-activity-dependent plasticity mechanism provides a compelling explanation for the formation and expansion of synaptic domains during development. Synapses that participate in coordinated local activity within a domain are reinforced, while out-of-sync synapses are weakened, leading to the progressive organization of dendrites into distinct functional domains.

\section{Discussion}

Our study provides the first in vivo mapping of individual functional synapses on developing cortical dendrites before eye opening, revealing that organized synaptic clustering emerges through spontaneous activity-dependent mechanisms during the second postnatal week. Three principal findings emerge from this work.

First, we demonstrate that functional synapse density on layer 2/3 pyramidal neuron dendrites increases dramatically between P8 and P13, with both spine and shaft synapses contributing. This timeline aligns with electron microscopy data showing peak synaptogenesis near P14 in mouse visual cortex \citep{lendvai2000experience} but extends those findings by revealing the functional organization of these synapses in vivo. The observation that approximately 20\% of synapses account for the majority of transmission events suggests that synaptic strength is highly heterogeneous from the earliest stages of circuit assembly.

Second, we identify a novel spatial organization principle: synapses do not appear uniformly along the dendrite but instead emerge in distinct domains of co-active inputs that expand during development. At early ages (P8--P10), these domains are small and confined, leaving large stretches of dendrite without functional inputs. By P12--P13, domains have expanded to cover nearly the entire dendritic arbor. This progression suggests that domain formation begins with the seeding of a few co-active synapses at specific locations, followed by the recruitment and stabilization of additional synapses through local plasticity mechanisms.

Third, and most critically, we show that local co-activity with neighboring synapses predicts synaptic fate: co-active synapses are potentiated and stabilized, while non-co-active synapses are depressed and potentially eliminated. This finding provides direct evidence for a Hebbian-like plasticity rule operating at the dendritic level \citep{hebb1949organization, markram1997regulation}, where the relevant ``postsynaptic'' signal is not the global somatic response but the local dendritic calcium level generated by nearby co-active inputs. This mechanism can explain how dendrites become organized into functionally distinct domains through a self-reinforcing process: once a few synapses begin firing together (driven by spontaneous network activity), they create a local calcium signal that promotes the potentiation of co-active neighbors and the elimination of non-co-active synapses.

Importantly, all of our recordings were performed before eye opening (P8--P13; eye opening occurs at P14 in mice), demonstrating that visual experience is not required for the initial establishment of synaptic clustering. Instead, the spontaneous network activity present during this period \citep{blankenship2010mechanisms, valeeva2019emergence} appears sufficient to drive domain formation. This spontaneous activity generates coordinated barrages that activate multiple synapses within local dendritic segments, providing the co-activity patterns needed for the plasticity mechanism we describe.

Our approach overcomes a major technical barrier that had prevented the mapping of shaft synapses in vivo. By including QX-314 in the intracellular solution to block action potentials and holding neurons at $-30$~mV to enhance NMDA receptor-dependent calcium influx, we could detect synaptic calcium transients at shaft locations that would otherwise be masked by back-propagating action potentials. This is critical for studying developing circuits, where shaft synapses are abundant before spine maturation \citep{kleindienst2011activity}.

Several limitations should be considered. First, our recordings were performed under light isoflurane anesthesia, which may attenuate the amplitude and frequency of spontaneous activity. However, prior work has validated that basic spontaneous activity patterns are preserved at the concentrations used. Second, the relatively small number of experiments per age precludes detailed statistical analysis of age trends within narrow time windows. Third, we cannot rule out the possibility that the patched neuron is not fully representative of the broader population.

In summary, our findings establish that functionally organized synaptic clustering on cortical dendrites is established by spontaneous activity before the onset of sensory experience, through a local co-activity-dependent plasticity mechanism that selectively potentiates synchronized synapses and eliminates non-synchronized ones. This developmental program creates the dendritic domain architecture that will later support the complex dendritic computations characteristic of mature cortical circuits.

\section{Conclusion}

This study reveals that functionally clustered synaptic inputs assemble into distinct dendritic domains on layer 2/3 pyramidal neurons in mouse primary visual cortex before eye opening. Using a combined in vivo two-photon calcium imaging and whole-cell patch-clamp approach that enables detection of both spine and shaft synapses, we mapped 354 synapses across 12 dendritic areas during the critical P8--P13 developmental window. Synapse density increased several-fold, domains of co-active synapses expanded from confined segments to near-complete dendritic coverage, and local co-activity predicted synaptic potentiation---demonstrating that spontaneous network activity drives the formation of organized synaptic architecture through a local plasticity mechanism, independent of sensory experience.

\bibliographystyle{plainnat}
\bibliography{references}

\end{document}
